\documentclass[11pt]{article}
\usepackage{graphicx}
\usepackage{times}
\usepackage{mysects}
%\usepackage{lucidbry}
\usepackage[active]{srcltx}
\newcommand{\figref}[1]{fig.\ \ref{#1}}
\newcommand{\bfig}{\begin{figure}}
\newcommand{\efig}{\end{figure}}
\newcommand{\bcen}{\begin{center}}
\newcommand{\ecen}{\end{center}}
\newcommand{\bali}{\begin{align}}
\newcommand{\eali}{\end{align}}
\newcommand{\beqn}{\begin{eqnarray}}
\newcommand{\eeqn}{\end{eqnarray}}
\newcommand{\beq}{\begin{equation}}
\newcommand{\eeq}{\end{equation}}
\newcommand{\benu}{\begin{enumerate}}
\newcommand{\eenu}{\end{enumerate}}
\newcommand{\bdes}{\begin{description}}
\newcommand{\edes}{\end{description}}

\newcommand{\km}{\, \textrm{km}}

\textwidth 6.2in \hoffset -0.5in
\textheight 9in
\voffset -0.7in  %  may be necessary depending on the latex implementation
%\renewcommand{\baselinestretch}{1.1}
\sloppy
%\pagestyle{myheadings}
\markright{\textit{Climate Prediction}}
%\markboth{\textit{Draft Pre-PreProposal}}{\textit{Climate Prediction}}
\begin{document}

\title{{Climate Predictability and  Variability in the Presence of
Anthropogenic Climate Change: \\
  A Proposal for a Princeton Climate Prediction Center}}
\author{\sc Prepared by G. K. Vallis \textit{et al.,}\thanks{G. K. Vallis
and J.~L.~Anderson, A.~J.~Broccoli, T.~Delworth, K. Dixon, I.~Held,
N-C.~Lau, S.~Griffies, J.~Mahlman, S.~G.~Philander, V.~Ramaswamy, A.
Rosati, J. Sarmiento, T.~Schneider, R.~Stouffer, R.~Toggweiler and
many others in the GFDL and Princeton University community.}
 \\ June 24 2000}
\date {\vskip -0.8 cm {~}}
\maketitle

\begin{narrower}
We propose the establishment of a Center in climate prediction,
involving both GFDL and Princeton University.  Our ultimate goal is
to develop a predictive capability for the climate system on
timescales from the seasonal-to-interannual, on which the initial
conditions are of paramount importance, to the multi-decadal, upon
which the signal from anthropogenic climate change may begin to
dominate. The activities of the Center will be designed to complement
and supplement existing research at GFDL and Princeton University in
such a way as to create a pre-eminent center for climate prediction,
fully integrated with NOAA's emerging Climate Services.

The activities of the Center will include the continual production
of ensembles of climate integrations for the next several decades,
eventually incorporating the present state of the climate system as
initial conditions. The integrations will be accompanied by
assessments of the confidence we can place in them, and thereby
assessments of the predictability of climate on interannual to
multi-decadal timescales. Our predictions will be continually tested
using the available observational networks, and the data and reports
will continually be made available to the public and the scientific
community.  The accompanying scientific activities will include
investigations into the mechanisms of climate variability on
interannual to multi-decadal timescales, including the variability
that has occurred in the past, and the interaction of natural
variability with anthropogenic climate change.


\bigskip
\hrule
\bigskip

\end{narrower}


\section{Our Goals} % and Products}

The simulation and understanding of natural variability on
multi-decadal scales intersects the issue of human-induced climate
change in a variety of ways. It begets the question of whether
predictions\footnote %
{On these long timescales, the word `projection' is often preferred
  over `prediction' because of the virtual impossibility of predicting
  global-scale human decisions. For simplicity we use the word
  prediction throughout this document, with the understanding that our
  predictive capability may be limited and that long-term
  anthropogenic climate change is not just a prediction in the
  usual initial-value-problem sense.}%
 of climate a decade or more in advance can be improved by
initializing and predicting the evolution of natural variability on
these time-scales. This question motivates our proposal.  Our goal is
the creation of a prediction/assessment capability that seamlessly
connects the seasonal-to-interannual timescale to both natural and
anthropogenic changes on longer time scales. Our explicit goal is to
provide predictions of climate and climate change over the next
50--100 years, incorporating estimates of all likely relevant future
forcing terms, and initializing the state of the climate system with
observed variables. The forcings will include an array of greenhouse
gases, with an interactive carbon cycle, as well as statistically
based estimates of the effects of future volcanic activity and
alteration of the land surface, for example by deforestation. The
initial states will include, in particular, the current oceanic state
as observed by both \textit{in situ} and remote instruments,
including a network of floats and satellites. Our experimental
predictions will ultimately be based on an ensemble of model
integrations, and so will be probabilistic in nature, with a
distribution of predicted states determined by the uncertainty in the
initial conditions, model errors and errors in model forcing,
filtered and amplified by the natural variability and
unpredictability of the climate system. The quantification of climate
predictability and variability will be a key element of the Center.

In summary, we seek to be the leading provider of accessible and
scientifically credible climate predictions in the US, in addition to
providing scientific research in support of more operational
predictions elsewhere. We envision that the information we will
provide will be invaluable to those performing impacts studies and
those involved in policy, in addition to the climate change and IPCC
communities. The Center will build on expertise at both GFDL and
Princeton University. It will interact with and adjoin the existing
research groups at GFDL and Princeton University, and it will
interface with observational and data analysis programs elsewhere
(e.g., PMEL, AOML, NODC). Overall, we intend and expect that the
Center will play an integral role in NOAA's emerging Climate
Services program.


\section{Our Product}

The primary product from such a Center will be a series of
experimental climate change predictions from the
seasonal-to-interannual to the multidecadal timescale, as well as
research into climate predictability and development of methods for
such predictions. As some of these predictions become more routine,
we envision that the techniques developed  will be passed on to more
operational centers within NOAA. Nevertheless, for the next decade
at least, predictions on the longer timescales may be performed in a
research environment side-by-side with research activities into the
nature of and limits to the predictability itself.

The probabilistic nature of climate change will necessitate that our
predictions primarily be made in the form of changes to the
probability distributions of climate. While the current
generation of models is not capable of making credible predictions of
climate change on the regional scale, it is expected that research
conducted in this Center will facilitate the future production of
such regional climate change forecasts, but the timescales on which
such forecasts will be credible is not as yet determined.

Our model results, and the interpretation of these results, will be
made available via the publication of periodic assessments, or
reports, and periodic briefings to the interested community. The
frequency of such reports will vary according to the timescale of
concern, perhaps varying from an annual to a five-year period.  In
addition, we propose to establish a state-of-the-art web site for the
dissemination of both the raw model data and the processed and
interpreted data. We envision this web site as comprehensive, and of
use to the general public, the private and public sectors, and the
scientific community. For the general public, our goal is for the
site to serve as the focal point in the U.S. for information on
potential climate change. This will encompass the scientific basis of
climate change, the characteristics and likelihood of potential
future changes, and possible impacts. Indeed, we would like this site
to be a catalyst for generating interest in the general public and
the media on these issues. Similarly, we envision that both private
and public sectors will benefit enormously from credible and
accessible assessments of the likelihood of future changes in (for
example) temperature and precipitation.  Energy and water resources
planning are but two obvious examples. And as our knowledge of
regional climate change improves, the benefits to these communities
become correspondingly greater.

For the scientific community the web site will serve as an easy
access to a wide variety of models and model-based products. A
specific effort will be made to provide easy access to model output
for the climate change impacts community. [Software developed at PMEL
(the `Live Access Server') will greatly help in this regard.]  Scientific communities (such as IRI and IPCC) will also
be able to access customizable subsets of model output for their own
assessments of its usefulness and impact. Finally, and importantly,
our scientific reports and publications in the refereed scientific
literature will provide an archival and high quality foundation for
all of the activities we describe.

The effort involved in all these activities and products will, of
course, be substantial. However, if realized properly, it will be a
valuable resource for NOAA and the nation. The following section
describe the scientific problems associated with such an endeavour,
and is followed by a summary of the proposed activities (see also
fig.\ 1).

\section{Climate Predictability}

The predictability of any system is normally defined in terms of how
well one can predict, or forecast, the future state of the system,
generally compared to some prior estimate that one can obtain
trivially, like persistence or climatology. The predictability of both
weather and climate is limited because they are chaotic dynamical
systems. Daily weather is essentially unpredictable beyond a couple of
weeks, because the information in the initial conditions is largely
lost and forecasts show little more skill than climatology.
However, the climate system as a whole may be much more predictable
than this, because many components (such as the ocean) vary much more
slowly, and these provide slowly varying boundary conditions for the
atmosphere. On a \emph{daily} timescale the variability of weather
is sufficiently large that a detailed knowledge of these boundary
conditions provides only a very small increase in skill: essentially
the signal to noise ratio is very small. However, this does not
necessarily hold on longer timescales; the signal to noise ratio
improves and the average weather may indeed be more predictable than
a climatological prior if the state of the ocean is known. Indeed on
seasonal-to-interannual timescales there is evidence that numerical
forecasts can be made that improve on climatology (Anderson
\textit{et al.}\ 1999). Much of the source of this skill arises from
predicting the conditions of the tropical ocean, and the effects
this has on the global atmosphere. Of course different physical and
dynamical processes dominate at the various timescales involved, and
the way in which predictability is judged consequently must differ.
It is useful to identify three distinct regimes:

\benu
\item The seasonal-to-interannual timescale, on which the forecast
  problem is largely an initial value problem. Skill resides in
  predicting the atmospheric response to low-latitude SST anomalies,
  and thus the prediction of El Ni\~no events is particularly
  important. Soil-moisture and surface albedo are also important
  components of initial conditions that affect seasonal skill.

\item A longer, multiyear-to-interdecadal, timescale on which there is
  significant natural variability, likely with origins in the
  ocean and perhaps involving coupled ocean-atmosphere modes in the
  tropics and/or extra-tropics. Such long-term variability may also
  arise from variability of frequency and intensity of El Ni\~no
  events.

\item The timescale at which the response to greenhouse gases
  dominates that of natural variability. This timescale is uncertain,
  for it depends both on the rate of emissions of greenhouse gases
  and aerosols, and the nature of the averaging process used to
  define the signal.

\eenu

Effort in the community has focused largely on the first and third of
these. For example, one of the goals of the Experimental Prediction
Group at GFDL is to understand predictability at the
seasonal-to-interannual timescale, and the IRI at Lamont-Doherty
concentrates on forecasts and their associated impacts at this
timescale. The scientific problem is largely an initial value one. In
contrast, in most global warming scenarios the precise initial state
of the climate system has not been regarded as of first-order
importance, and for centennial timescales this may in fact be
appropriate. In contrast to either of these, for the decadal timescale
\textit{both} natural variability (and therefore the initial
conditions of the climate system) \emph{and} the forced response to
increasing greenhouse gases can be expected to be important. It is the
interaction of the initial value problem and the forced response that
is central to our proposed activities.  An illustration of the
importance of this emerges from the simulations of Delworth and
Knutson (2000). They performed an ensemble of five integrations over
the 20th century, using a climate model that differed only in its
initial conditions. In all five integrations there is significant
natural variability, and one member of the ensemble provides a
remarkably good match to the observed global temperature record
(\figref{delworth1} and \figref{delworth2}). The implication is
that \emph{both} natural variability and the forced response to
increasing concentrations of greenhouse gases are important on
such timescales.


\subsection{Causes of unpredictability in the climate system}

As in many predictability problems, lack of skill is
ultimately caused by:
\benu
 \item Imperfect models, including imperfect knowledge of forcing and
boundary conditions.
 \item Imperfect knowledge of initial conditions.
\eenu %
The errors from these sources are amplified by internal instabilities
and feedbacks in the dynamics of the system over which we have no
control.  From a practical standpoint then, our overall effort will
involve continually improving the numerical models we use to predict
climate and initializing our models with the best possible data. More
specifically, we expect that difficulties in climate prediction will
arise from our lack of understanding or ability in the following
areas:

\benu

\item In our ability to initialize the initial state of our climate
  models with a sufficiently accurate state of the ocean
  and sea-ice, and sufficiently accurate land-surface conditions.
\item \label{idec} In understanding the nature of decadal variability
  in the ocean--atmosphere--land--ice system. Specifically, we must
  better understand the nature and causes of various decadal-scale
  oscillations in the climate system, which in turn requires:
\benu
\item Understanding and modelling the mechanisms determining the
  evolution of the state of the ocean (e.g., SST anomalies) and sea ice.
\item Understanding and modelling the effects of such SST and sea-ice
  anomalies on the state of the atmosphere.
 \eenu
  Furthermore, because we cannot realistically expect that our models
 and data assimilation will ever be perfect, we must quantify the
 nature of climate variability that has occurred in the past, from all
 possible causes, in order to gain an appreciation for what might
 occur in the future.
\item In understanding and predicting the likely future course of the
  radiative forcing functions of the atmosphere, including greenhouse
  gases, sulfur and aerosols and, conceivably, volcanic events. The
  variability of both the terrestrial and marine components of the
  carbon cycle are also major factors in this.
\item In understanding which components of the climate system are
      predictable, and quantifying the level of predictability we
      might attain. This in turn entails understanding the complex
      feedbacks that determine the nature of the response
      of the climate system to natural and anthropogenice forcing.
\eenu

These issues will now each be discussed in more detail.

\section{Obtaining an Initial State: Ocean Data Assimilation}
On timescales beyond a month, the signal in the climate system will
largely arise from the slowly varying components --- primarily the
ocean and sea ice, and perhaps the land surface.  The dynamics of the
ocean is known to contain timescales of decades (the wind-driven
circulation) to centuries and even millennia (the abyssal
circulation), and thus to capture any climate predictability we must
initialize our climate models with the current state of the ocean. We
do not know how long the influence of such initial conditions will be
felt. Nor do we know well how much inertia resides in the land
surface, in particular soil moisture content and albedo, but these
too must ultimately become part of the initial conditions.

% Although some studies with idealized atmospheric
% general circulation models indicate that the atmosphere alone can
% exhibit a red spectrum (e.g., James and James), initializing the
% atmospheric state with observations is unlikely to lead to increased
% predictability on climatic timescales.

A number of ocean data sets exist, including the well-known Levitus
data set and data emerging from the WOCE program. We also propose
to take advantage of new data sets that will be emerging over the next
few years, in particular the Argo float data set, and SST and
sea-surface height (SSH) data from satellites. The Argo program (Argo
1998, 2000) is a follow-on to the WOCE hydrographic, float and
expendable bathythermograph sampling networks, but uses an array of
relatively cheap, profiling PALACE-like floats (e.g., Davis 1998).
The goal of the program is to have a global array of such floats, with
one float profiling to about 2000 m depth about every 10 days in every
3-degree square of the ocean, providing measurements of temperature,
salinity, pressure and (from the float trajectories) a velocity field
at the reference depth. All told, about 80,000 T/S profiles will be
taken per year, sufficient to provide a good picture of the
large-scale baroclinic structure of the upper ocean, including some
information below the main thermocline. The floats do not sample the
abyss, but this is both much more uniform than the upper ocean, and
the typical timescales on which it varies are for the most part much
longer than the decadal timescales that concern us here.  Nor is the
density of the floats sufficient to resolve mesoscale eddies.

In contrast to floats, satellite altimeters (such as TOPEX/POSEIDON,
and the next generation Jason altimeter) provide relatively high
resolution information about sea-surface height, including the
oceanic eddy field. The information is complementary to that of the
Argo program, because floats provide the subsurface vertical
structure at relatively coarse horizontal resolution whereas
altimeter data provides a higher resolution (in time and, especially,
space) picture of sea-level variability. Sea surface temperature is
also routinely available from satellites.

The cost of the Argo program is in excess of \$40 M, and that of the
Jason and other satellite missions much higher still. Assimilating
this data into ocean models, and using these models for climate
prediction, greatly increases the value of the data, and will be an
integral part of the proposed project. We expect that this aspect of
the proposed work will be conducted in collaboration with other NOAA
laboratories, in particular PMEL, AOML and NODC, and as part of the
NOAA Ocean Data Observations and Services initiative. Indeed, our
proposed work provides a natural and important context for this
activity. R.~Molinari of AOML is on the Argo Science Team and
has expressed interest in collaborating with GFDL in this endeavour.
Sydney Levitus of NODC, perhaps the leading provider of oceanic data
sets in the country, has expressed a similar interest.


\section{Climate Variability on the Decadal to Multidecadal Timescale}

Over the past several years a bewildering variety of oscillations and
mid-to-high-latitude hemispheric, decadal-scale phenomena have
emerged. These include the Pacific Decadal Oscillation (PDO), the
North Atlantic Oscillation (NAO), the Northern-Hemisphere Annular Mode
(NAM), the Arctic Oscillation (AO), the Antarctic Oscillation (AAO) and
the Antarctic Circumpolar Wave (ACW). See, among many others, Walker
and Bliss (1932), Bjerknes (1964) and Thompson and Wallace (2000).
These modes (not all independent) are in some contrast with their
better-known low-latitude counterpart, the Southern Oscillation, part
of the El Ni\~no-Southern Oscillation phenomenon. Typically, the mid-
and high-latitude oscillations have a much longer period and a less clear
signal in terms of, say, a surface pressure index than ENSO. The modes
are nevertheless potentially very important for climate variability
and predictability --- indeed Wallace (2000) comments that `This
phenomenon [the NAO/NAM] rivals the El Ni\~no Southern Oscillation
(ENSO) in terms of its significance for understanding global climate
variability and trends.'

The existence of any putative quasi-regular phenomena on a decadal
timescale would indicate that the phenomena might in principle be
predictable, and some recent numerical work is encouraging in this
regard. For example, recent work by Griffies and Bryan (1997) suggests
that the state of the interior North Atlantic ocean \emph{is}
predictable for a decade or longer.  Furthermore, Rodwell \emph{et
al.}\ (1999) find that they are able to simulate trends in the NAO
(an atmospheric index) over the past few decades by imposing the sea
surface temperature and allowing the model atmosphere to respond.
Taken together, these results might imply that there \emph{is} useful
predictability in the climate system on the decadal timescale, and
that our numerical models may be close to being able to capture this
predictability. Such conclusions will only hold if the evolution of oceanic
SST anomalies is largely determined by the ocean's own internal
dynamics. The fully-coupled ocean-atmosphere system may be less
predictable, because, for example, the evolution of SST anomalies may
be largely determined by stochastic and essentially unpredictable
atmospheric processes (Bretherton and Battisti 2000).  Also, the trend
in the AO may be due to a change in atmospheric dynamics associated
with increased greenhouse gas concentrations (Shindell \textit{et
  al.}\ 1999). Certainly, it is not clear that oscillations such as
the NAO are any more than red-noise (Wunsch 1999), and the observed
trends may simply be a natural fluctuation in that noise. In any case,
the consequences of changes in such oscillations would be profound,
and this alone dictates that a considerable effort should be
undertaken to better understand and predict them. For example,
Delworth and Dixon (2000) show that the recent upward trend in the AO
index, leading to an increase in the westerly winds over the mid and
high latitudes of the North Atlantic, could substantially delay the
weakening of the thermohaline circulation that is projected as a
consequence of increased high latitude precipitation in a warmer
world.

Another facet of predictability is the possible decadal variability of
ENSO itself. Over the past century certain decadal long periods have
been virtually bereft of El Ni\~no events (for example the 1930s),
whereas in other decades their frequency and amplitude is abnormally
high.  Recent work indicates that variations in the thermocline depth,
possibly caused by decadal scale oscillations in mid-latitudes, may be
the underlying reason. With global warming however, the abyssal ocean
is expected to warm more slowly than the upper ocean, the depth of the
main thermocline may change as a consequence, and there may be
potentially significant associated changes in El Ni\~no frequency and
intensity. This example further underscores the importance of the
interaction of natural variability with anthropogenic climate change,
which lies close to the heart of our proposed work.

Unfortunately, the current generation of fully-coupled
ocean-atmosphere climate models are currently unable to reproduce
trends in the mid-latitude oscillations themselves, nor can they
properly produce El Ni\~no events with the right magnitude and
frequency. We need both better models and a better understanding of
the dynamics of such oscillations in order to have confidence in our
models' ability to correctly capture and predict them; a similar plea
is made also by Wallace \textit{(op cit).} In addition to a
continuing effort at GFDL to improve the quality of our GCMs, and
the use of such models to study ocean-atmosphere interaction,
research activities in the Center must involve modelling and
theoretical studies that allow feedbacks and oscillations to be
studied in a conceptually simpler framework (e.g., Lau and Nath
1994, Saravanan \textit{et al.}\ 2000). One potentially large source
of predictability decay lies in the effects of eddies on the
large-scale circulation of the ocean: an upscale cascade of energy
from the eddy scale might provide an additional, perhaps
unpredictable, source of variability to the large-scale
circulation.  The present (and next) generation of computers are
insufficiently powerful to incorporate this effect in the global
models to be used for climate predictions, so such effects must be
studied using more idealized models and included parametrically in
the global models, at the same time as we prepare for the future use
of fully eddy-resolving ocean models in climate simulations.

\section{Quantifying Natural Variability and Extreme Events}

As we have stated, climate predictions will be expressed in terms of
probabilities, using ensembles of integrations differing in initial
conditions and other parameters. All deterministic methods can be
developed for climate prediction are certainly subject to
uncertainties, and even the level of uncertainty is uncertain.  Thus,
although we will estimate a probability density function (PDF) for the
future climate state from the ensemble, we would like to have an
independent check on its width.

Furthermore, many of the most significant climate impacts may involve
changes in high impact events such as droughts, floods, hurricanes and so
forth, so that the type of climate projections with the most economic
value will likely involve the probability of such events. The
possibility of local or global \emph{rapid climate change} is another
issue that demands similar attention. Prediction of such phenomena is
one of the hardest tasks to ask of climate models, for it involves
understanding how subtle changes in the global circulation might lead
to changes in the frequency or intensity of much smaller scale
phenomena, which are typically not well resolved by the global models.

Two approaches will be used to address this issue; one is primarily
model-based, the other involves extensive use of historical data.  In
the first, we use high resolution, regional models to examine how
smaller scale phenomena might change in future climates. For example,
in a prototypical study, Knutson \textit{et al.} (1998) examined how
hurricane intensity might change in a warmer world, using model data
from a global warming simulation to provide boundary conditions for a
high resolution, hurricane model.  They found hurricane intensity
typically increased in a warmer climate.  Similar techniques can be
used to examine how the behavior of other small-scale phenomena might
be modified in response to a changed large-scale circulation. The
development of the powerful Flexible Modelling System (FMS) at GFDL
will allow changes in model configuration and resolution to be made
relatively easily, greatly facilitating the way such studies can
be performed.

But even if our models were adequate to accurately simulate such
small-scale phenomena, practical limitations will likely limit the
number of ensemble members to a value well below that required to
adequately estimate the tails of the PDF of the climate state.
Thus, for our second approach, we turn to the historical record to
attempt to quantify the natural variability of the climate system. One
particular methodology assumes that changes in variability due to
anthropogenic forcing will be small, and simply shifts the PDF based
on changes in central value.  To use this approach, very good
quantitative estimates of the PDFs of the climatic quantities of
interest are required.  Considerable data are available for estimating
such PDFs. During the 20th century, the instrumental climatic data
available over North America and Europe are particularly good,
especially in the latter half of the period, but the duration of the
data is not sufficient to allow an accurate assessment of the
probability of unlikely events; that is, to determine the tails of the
PDF.  This problem is exacerbated by the high likelihood that
anthropogenic trends are present during this period, which makes it
difficult to isolate the low-frequency internal variability of the
climate system.

To address this issue, development of paleoclimatic proxies with
annual resolution have allowed very long (up to 1000 years) climatic
time series to be developed for certain sites.  There are also a few
long instrumental records extending back to the late 18th and early
19th centuries. Such records, whether proxy-derived or directly
measured, have the promise of delivering important information about
climate variability. Unfortunately, the spatial coverage of these
long records is very limited.  Innovative techniques (e.g., Mann and
Bradley 1999) have allowed the reconstruction of large-scale
temperature patterns, from climatic proxies, but the fidelity of such
reconstructions is probably limited to relatively large scales. Some
of the most important climate impacts result from smaller-scale
anomalies, so a strategy must be developed for estimating the
frequencies of such events at scales unresolved by the data.

Long climate integrations can aid in this effort, and part of our
proposed activities will involve creating an ensemble of very long
integrations of a climate model, for example for the last millennium,
using the most accurate forcing functions possible as developed by the
Atmospheric Processes Group at GFDL. This includes the effects of
variations in greenhouse gases, aerosol content, volcanic activity and
solar variations over the last three hundred years, with less
accurate estimates of these forcings over the last one thousand years.
Then, if adequately calibrated using existing paleoclimatic and long
instrumental records, estimates of climate variability from models can
be used to fill the spatial gaps between such long records. Because,
they can be run for periods much longer than our best records, such
model integrations can be used to better estimate the tails of the
climate PDFs, and thus the return periods of infrequent events and the
likelihood of such events occurring in the future.


\section{Radiative Forcing and the Carbon Cycle}

On the multi-decadal timescale, determination of the levels and the
interaction of the external forcing agents becomes a key issue in
quantification of climate change. These agents include the major well-mixed
greenhouse gases such as CO$_2$ and methane, as well as
sulfur and particulate carbon. On the multi-decadal to centennial
timescales, as the climate system loses its memory of the initial
conditions our predictions become increasingly dependent on various
emissions scenarios. The human-influenced sources and sinks of
greenhouse gases and aerosols become nearly impossible to predict,
but this merely underlines the importance of understanding their
effects in order to be able to make rational policy decisions.
Indeed, future projections of greenhouse gases and aerosols are
needed for at least three reasons: (i) to help quantify anthropogenic
climate change; (ii) because of the potential response of terrestrial
and marine biota to climate change; (iii) because the development of
sensible emission and mitigation strategies is dependent on making
such projections as quantitative as possible.

Carbon uptake is one significant source of uncertainty in global warming
scenarios: there is considerable uncertainty about how much this
uptake might vary naturally on seasonal-to-interannual to decadal
timescales, and how this may be affected by increasing
industrialization. In order improve such projections, we need to be
able to answer the following fundamental
questions: %
\benu
 \item  Where and how large are the main current terrestrial and oceanic
sources and sinks of CO$_2$?
 \item How are these likely to change in the future,
both as a result of natural variability and anthropogenic climate
change?
\eenu
To answer these requires in turn an understanding of how the carbon
cycle and the climate system interact and function together, and this
will be an important goal of the Center. Considerable progress
has already been made toward this by the Carbon Modelling Consortium
(CMC) at Princeton University, but understanding how the variability of
the biogeochemical processes is linked to that of the physical
processes remains a key problem in predicting climate change. The
interaction of other elements is even less understood,
especially non well-mixed greenhouse gases and particulate matter.


Thus far, most of the assessments of the oceanic uptake of
anthropogenic CO$_2$ have assumed that the biological system is not
altered by climate change. Thus the chemical solubility of the ocean,
but not potential changes in biological activity, have been included in
the models. There is increasing evidence that this is unrealistic and
overconstrains the models (Denman \textit{et al.}\ 1996). The response
of such models both to natural variability and to anthropogenic climate
change may thus be too limited, potentially resulting in an
poor estimate of the variability of future levels of CO$_2$.

Similar issues arise in the terrestrial biosphere, where the feedback
between climate and vegetation could lead to human activity having
very untoward effects. In the most extreme case, deforestation could
lead to a dry tropics, which would then significantly feed back onto
the climate system as a whole. Newly developed models of the
terrestrial biosphere, based on appropriately scaled-up models of
communities and ecosystems, can be used to ascertain whether such
scenarios are at all tenable. However, the complete incorporation of
both terrestrial and marine ecosystem models into the GCMs is really
necessary to understand the full range of feedbacks between the
physical and biogeochemical systems, and to obtain accurate
predictions of likely future CO$_2$ levels.




\section{Detection, Attribution and Prediction}

The issues involved in `detection and attribution' play a role not
only in how climate has changed thus far in response to increasing
greenhouse gases, but also in quantifying future changes and the
plausibility of such changes. In practice, detection usually means
that there is a significantly greater degree of resemblance between
observations and a climate simulation with anthropogenic forcings than
between observations and a simulation without anthropogenic forcings.
The degree of resemblance is usually measured by a regression
coefficient of observations onto a model-simulated climate change
signal, say the average difference between a simulated reference
climate and the simulated response of the climate system with
increased concentrations of greenhouse gases and aerosols (e.g.,
Santer \textit{et al.\ }1995; Hegerl \textit{et al.}\ 1996).
% Stott and Tett 1998).
This type of analysis is sometimes known as \emph{optimal
  fingerprinting}. Attribution goes further and attempts to associate
signals with particular causes, again via the use of model simulations
and pattern matching (e.g., Hasselmann 1997). An irreducible physical
difficulty in all these procedures comes from the natural variability
in both the climate and the simulations.  Natural variability and
anthropogenic climate change are frustratingly intertwined,
and no matter what statistical tests are employed, the natural
variability of the system must be properly represented in the models
or the answers may be misleading.

Similar issues recur in the predictability problem, because of the
need to understand the degree to which the climate response is due to
natural variability or to a forced response to greenhouse gases. In
the standard problem (as it has been posed in the past) one simply
wishes to distinguish the forced response due to greenhouse gases
from that due to natural variability (`noise'). We propose to go one
step further, and to understand what aspects of the \emph{natural
variability} of the climate system are predictable, and can thus be
regarded as signal rather than noise. There are two classes of
questions we wish to address: %
\benu
 \item Which aspects of
the climate system are most predictable simply in the context of
global warming? That is, which aspects of the future climate system
robustly differ from a climatology with no increased greenhouse
gases?
 \item  Which aspects of the climate system can be better
predicted by initializing the model from data? Conversely, what
aspects of the initial field is the future evolution of the climate
most sensitive to?
\eenu

Key to answering these important questions will be performing multiple
\emph{ensembles} of climate integrations, initialized with differing
oceanic and atmospheric states, to isolate the role of natural
variability and the initial conditions.  Given such an ensemble,
various statistical tests may be employed to help answer these
questions (e.g., Anderson and Stern 1996; Schneider and Griffies
1999). For example, we can expect the SST of the North Atlantic to be
predictable, for some time period, if the initial data are given.
What other fields, especially atmospheric ones, might be predictable
in this sense? The quantification of such issues, indeed the
quantification of climate predictability, will be a central theme of
our Center.  The combination of carefully designed ensemble
experiments, assimilating data to determine the initial conditions,
and new statistical methods, is a novel and powerful approach to this
problem.



\section{Research Activities: an Overview}

Let us now summarize some of the research that we propose, including
the numerical experiments that will be done, the subsequent analysis,
and the ongoing basic research.

We expect to be continually performing multiple ensembles of
experiments to examine the evolution and predictability of the
climate system. For example, we will perform a set integrations
which differ in their initial conditions in the ocean and sea ice,
since these are the slowly varying components of the system. We
intend that the ocean state should incorporate assimilated
observations, with other ensemble members being perturbations around
that state.  Since practical issues will force the ensemble to be
small, it will be important to carefully choose the ensemble both so
that the initial error is representative of our lack of knowledge of
the initial conditions and so that the oceanic error growth is not
artificially small. Effectively, the PDF of the data error must be
assimilated in a manner similar to that of the data itself.

In addition, we will form an ensemble of integrations with differing
\textit{atmospheric} initial conditions but identical oceanic
initial conditions. If there is significant climate predictability,
the ensemble members that differ only in initial atmospheric state
should track each other better than members with differing oceanic
initial conditions. Our statistical predictability analyses will
quantitatively assess this, determine which components of the
climate system might be predictable, and which aspects of the
initial state the evolution of the system is most sensitive to.  We
again emphasize that for such experiments to be meaningful, climate
models that plausibly simulate the natural variability of the
climate system must be employed. This will entail continual model
development in conjunction with existing activities at GFDL, and
continual testing of the predictions of the models against suites of
oceanic and atmospheric observations.

We are now actively engaged in the planning of our first set of
experiments. We will take three different oceanic initial conditions
corresponding to calendar year 2000 from three different transient
greenhouse gas integrations. Thus, we have three different but
plausible oceanic initial states.  Each case will then be run for 50
years, each with three different atmospheric initial conditions,
resulting in nine total experiments. A much smaller ensemble will be
continued out to 100 years. We envision that this kind of activity
will go on continually throughout the lifetime of the Center, using
increasingly realistic initial states and increasingly more
realistic models and forcing functions, and would be the basis for
the assessments produced.

In parallel and coordinated with this activity will be research into
the causes and nature of climate variability and the mechanisms that
limit its predictability. This will entail research into the various
feedback mechanisms that amplify and diminish the anthropogenic
forcings, and research into the forcings themselves. This will
include research into the variability of the carbon cycle and the
terrestrial and marine uptake of carbon dioxide, research into the
nature of oceanic variability and its influence on the atmosphere,
and research into the nature and possibility of extreme climate
events and rapid climate change. Study of the historical record in
conjunction with long model integration will be important in
quantifying natural variability and will serve as an important check
on our models.


\section{Structure and Composition}

%\begin{table}
%\caption {Participants in the Princeton Climate Center} \label{tab1}
%\begin{tabular}{ll}
%\\ \hline \hline \\
%At Princeton/GFDL & Geoffrey K. Vallis \\
%At GFDL & Isaac Held    \\
%        & Tom Delworth  \\
%        & Tony Broccoli \\
%        & Robbie Toggweiler  \\ \\
%At Princeton University & Jorge Sarmiento \\
%                        & George Philander \\
%                    %   & Steve Pacala  \\
%\\ Other Institutions \\
%NOAA/AOML   &  Robert Molinari \\
%NOAA/NODC   &  Sydney Levitus  \\
%\hline
%\hline
%\end{tabular}
%\end{table}

The proposed Center will be composed of both GFDL and Princeton University
scientists, and collaborators from other institutions.
%Some of the main participants are listed in the table. \ref{tab1}.
Although a very strong foundation for this activity already exists at
Princeton University and GFDL, work of the scope we have proposed
will require a major enhancement of the modelling and analysis groups
that already exist. We expect that significant progress could be made
toward our ultimate goals by the addition of about 20 scientists and
support staff, who would bring important new talents to the difficult
problems we have outlined. In particular, scientific teams would need
to be formed or greatly strengthened in the following areas:

\benu
\item Outreach to and interface with the impacts community.
\item Ocean Data Assimilation.
\item Model/data diagnostics and dynamics, including
      decadal/centennial dynamics and analysis of the historical
      record.
\item Model development and radiative forcing,
      including improved schemes for aerosols and marine and
      terrestrial biogeochemical cycles.
\eenu

Centers similar to the one we propose do exist elsewhere in the
world: for example, the Hadley Centre for Climate Prediction and
Research is the main center in the UK for such activities. The Hadley
Centre is a little larger than the entirety of GFDL, and its
portfolio explicitly includes a mandate to simulate the present
climate and understand its natural variability, to understand the
factors controlling climate change, and to predict global and
regional climate change up to the end of the 21st century. The
Max-Planck-Institute f\"ur Meterologie in Hamburg also investigates
the nature of interannual and decadal climate variability and the
predictability of climate variability on interannual to decadal time
scales. All of these efforts are very substantial indeed.
Nevertheless, because of the combined expertise of GFDL and Princeton
University and the strategy we have outlined here, we believe that we
are in in a position to make a unique and important contribution and
become the leader in this field.

We expect that the policy aspects of climate and climate change will
become increasingly important in years to come; the high-level
discussions stemming from the Kyoto Protocol are one indication of
this. At Princeton University, the Princeton Environmental Institute
(PEI) explicitly seeks to make connections between science and
policy, and we expect the Center to have considerable interaction
with PEI. Indeed, PEI has explicitly stated its willingness and
desire to foster such interaction, and ultimately to help build a
policy side to the Center. This will also involve the Woodrow Wilson
School of Public and International Affairs, whose faculty are already
concerned with a variety of policy issues related to the environment
and energy. The eventual establishment of a policy component within
the Center will enable scientists and policymakers to interact much
more closely than has been possible in the past, and consequently
the policy implications of climate change to be explored in much
more depth.

For a Center to embrace all of the activities we have outlined in
this prospectus will require the development of a new level of
infrastructure between GFDL and the University. One possible avenue
for this is the development of a Joint Institute, but the development
of any such infrastructure must be done carefully and thoughtfully.

% One possible avenue
% for this is the formation of a Joint Institute between GFDL and
% Princeton University. Prior to that, the Center can be funded by
% separate awards to Princeton University and GFDL.


%\subsection*{Acknowledgements}
%\bigskip \hrule \bigskip
%This document was prepared by G. K. Vallis,
%J.~L.~Anderson, A.~J.~Broccoli, T.~Delworth, K. Dixon, I.~Held,
%N-C.~Lau, S.~Griffies, J.~Mahlman, S.~G.~Philander, V.~Ramaswamy, A.
%Rosati, J. Sarmiento, T.~Schneider, R.~Stouffer, R.~Toggweiler and
%many others in the GFDL and Princeton University community.

\newcommand{\jgr}[1]{\textit{J. Geophys.\ Res.,} \textbf{#1}}
\newcommand{\jc}[1]{\textit{J. Climate,} \textbf{#1}}
\newcommand{\qjrms}[1]{\textit{Quart. J. Roy.\ Meteor.\ Soc.,} \textbf{#1}}
\newcommand{\grl}[1]{\textit{Geophys.\ Res.\ Lett.,} \textbf{#1}}

\newpage

\section*{References}
\bdes

\item Anderson, J. and W. F. Stern, 1996. Evaluating the potential
  predictive utility of ensemble forecasts. \jc {9,} 260--269.

\item Anderson, J. H. van den Dool, A. Barnston, W. Chen, W. Stern and
  J. Ploshay 1999. Present-day capabilities of numerical and
  statistical models for atmospheric extratropical seasonal simulation
  and prediction.  \textit{Bull. Am. Meteor. Soc.,} {\bf 80,}
  1349--1361.

\item Argo Science Team, 1998. On the design and implementation of
  Argo: an initial plan for a global array of profiling floats.
  International CLIVAR Project Office Report 21, GODAE Report 5. GODAE
  International Project Office, Melbourne, Australia, 32pp.

\item Argo Science Team, 2000. Report of the Argo Science Team Meeting
  (AST 2) Southampton, March 2000. 35pp.

%\item Barnett, T.~P., and M.~E. Schlesinger, 1987.
%Detecting changes in global climate induced by greenhouse gases.
%{\em J. Geophys. Res.\/}, {\bf 92D}, 147772--14780.

\item Bjerknes, J., 1964. Atlantic air-sea interaction. \textit{Adv.
    Geophys.,} \textbf{10,} 1--82.

\item Bretherton, C. S. and D. S. Battisti, 2000. An interpretation of
  the results from atmospheric general circulation models forced by
  the time history of the observed sea surface temperature
  distribution. \grl {} (in press).

\item Davis, R. E., 1998. Preliminary results from directly measured
  middepth circulation in the tropical and south Pacific. \jgr {103}
  24619--24639.

\item Delworth, T. and K. Dixon. 2000. Implications of the recent
  trend in the Arctic Oscillation of ocean circulation and climate
  change. \jc{} (in press).

\item Delworth, T. and T.  Knutson. 2000. Simulation of early 20th
  century global warming. \textit{Science,} \textbf{287,} 2246--2250.

\item Denman, K., E. Hofmann and H. Marchant, 1996. Marine biotic
  responses to environmental change and feedbacks to climate. In In
  \textit{Climate Change 1995: The Science of Climate Change},
  Contribution of Working Group 1 to the Second Assessment Report of
  the IPPC, ed. J. T. Houghton \textit{et al.,} pp.\ 193--227.
  Cambridge University Press.

\item Hasselmann, K., 1997.  Multi-pattern fingerprint method for
  detection and attribution of climate change.  {\em Climate Dyn.},
  {\bf 13}, 601--611.

%\item Hegerl, G.~C., K.~Hasselmann, U.~Cubasch, J.~F.~B. Mitchell,
% E.~Roeckner,  R.~Voss, and J.~Waszkewitz, 1997.
%Multi-fingerprint detection and attribution analysis of
%greenhouse gas, greenhouse gas-plus-aerosol and solar forced climate
% change.
%{\em Climate Dyn.}, {\bf 13}, 613--634.

%\item Hegerl, G.~C., and G.~R. North, 1997.
%Comparison of statistically optimal approaches to detecting
%  anthropogenic climate change.
%{\em J. Climate\/}, {\bf 10,} 1125--1133.

\item Hegerl, G.~C., H.~{von Storch}, K.~Hasselmann, B.~D. Santer,
  U.~Cubasch, and P.~D. Jones, 1996.  Detecting greenhouse-gas-induced
  climate change with an optimal fingerprint method.  {\em J.
    Climate\/}, {\bf 9}, 2281--2306.

\item Griffies, S. M. and K. Bryan, 1997. Predictability of North
  Atlantic multidecadal variability. \textit{Science,} \textbf{275,}
  181--184.

%\item Huck, T. and Vallis, G. K. 2000. Use of a tangent-linear model
%for studying the variability and predictability of an ocean model.
%Submitted to \textit{Tellus.}

%\item James, I and James

%\item Knutson, T.~R., T.~L. Delworth, K.~W. Dixon, and R.~J. Stouffer,
%  1999.  Model assessment of regional surface temperature trends
%  (1949--97).

\item Lau, N-C and M. J. Nath 1994. A modeling study of the relative
  roles of tropical and extratropical SST anomalies in the variability
  of the global atmosphere-ocean system. \jc {7,} 1184--1207.

\item Rodwell, M. J., D. P Rowell and C. K. Folland, 1999. Oceanic
  forcing of the wintertime North Atlantic Oscillation and European
  climate. \textit{Nature,} \textbf{398,} 320--323.

%\item McCalpin

%\item Meacham

%\item Mann, M. E., R. S. Bradley and M. K. Hughes, 1999. Global-scale
%temperature patterns and climate forcing over the past six centuries.
%\textit{Nature,}  \textbf{392,} 779--787.

\item Mann, M. E., and R. S. Bradley, 1999.  Northern Hemisphere
  temperatures during the past millenium: inferences, uncertainties,
  and limitations.  \grl {26,} 759-762.

\item Santer, B. B., T. M. L. Wigley, T. P. Barnett and E. Anyamba,
  1996. Detection of climate change and attribution of causes. In
  \textit{Climate Change 1995: The Science of Climate Change},
  Contribution of Working Group 1 to the Second Assessment Report of
  the IPPC, ed. J. T. Houghton \textit{et al.,} pp.\ 407--443.
  Cambridge University Press.

% \item Saravanan

\item Saravanan, R., G. Danabasoglu, S. C. Doney, J. C. McWilliams,
    2000. Decadal Variability and Predictability in the Midlatitude
    Ocean�Atmosphere System. \jc {13,} 1073�1097.

\item Schneider T. and S. M. Griffies, 1999. A conceptual framework
  for predictability studies. \jc {12,} 3133--3155.

\item Shindell, D.T., Miller, R.L., Schmidt, G.A., and L. Pandolfo,
  1999. Simulation of recent northern hemisphere winter climate trends
  by greenhouse-gas forcing.  \textit{Nature,} \textbf{399,} 452-455.

\item Stern, W. F. and K. Miyakoda, 1995. The feasibility of seasonal
  forecasts inferred from multiple GCM simulations. \jc {8,}
  1071--1085.

% \item Stott, P.~A., and S.~F.~B. Tett, 1998.
% Scale-dependent detection of climate change.
% {\em J. Climate\/}, {\bf 11}, 3282--3294.

%\item Tett, S. F.~B., P.~A. Stott, M.~R. Allen, W.~J. Ingram,
%and J.~F.~B. Mitchell,  1999.
%Causes of twentieth-century temperature change near the earth's
%  surface.
%{\em Nature,} {\bf 399,} 569--572.

\item Thompson, D. W. J. and J. M. Wallace, 2000. Annular modes in the
  extratropical circulation. Part I: Month-to-month variability.  \jc
  {13,} 1000--1016.

%\item Thompson, D. W. J., J. M. Wallace and G. Hegerl, 2000. Annular
%modes in the extratropical circulation. Part II: Trends.
%\jc {13,} 1018--1036.

\item Wallace, J. M. 2000. North Atlantic Oscillation/annular mode;
  Two paradigms --- one phenomenon. \qjrms {126,} 791--805.

\item Walker, G. T. and E. W. Bliss, 1932. World Weather V.
  \textit{Mem. R. Meteorol.\ Soc.} \textbf{4,} 53--83.

\item Wunsch, Carl, 1999. The interpretation of short climate
records, with comments on the North Atlantic and Southern Oscillations.
 \textit{Bull. Am. Meteor. Soc.} \textbf{80,} 245-256.
\edes


%\end{document}

\begin{figure}
\begin{center}
\includegraphics[height=4.0in]{../figs/chart3.eps}
\vskip 1in \caption {Schematic of some of the main activities in the
proposed Center.}
\label{chart1}
\end{center}
\end{figure}


\begin{figure}
\begin{center}
\includegraphics[height=3.0in]{../figs/fig1.eps}
\caption {\small (a) Time series of global mean surface temperature from
observations (heavy black line) and the five  model experiments
(colored lines). (b) Same as (a), except that only one of the model
timeseries (`ensemble member 3') is shown (red line).}
\label{delworth1}
\end{center}
\end{figure}


\begin{figure}
\begin{center}
\includegraphics[height=4.0in]{../figs/fig2.eps}
\caption {\small Zonal mean anomalies of surface temperature for the
observations (upper panel) and for ensemble member 3 of
\figref{delworth1}, subject to a 10-year low-pass filter.}
\label{delworth2}
\end{center}
\end{figure}

%\begin{figure}
%\begin{center}
%\includegraphics[height=3.0in]{../figs/fig3.eps}
%\caption {\small Delworth fig 3.}
%\label{fig4}
%\end{center}
%\end{figure}

\end{document}
